\section{When statistical and local-query costs multiply}\label{sec:composition}
Direct sums preserve spectra and sparsity, but do not automatically
multiply query lower bounds. The missing quantity is the relative
variation of a \emph{single positive form}, rather than an off-diagonal
entry obtained by subtraction. We prove a distributional product that
avoids pointwise two-level promises, and then quantify two construction
classes that cannot supply the stronger condition-dependent product.

\subsection{Averaging and contrast}
Suppose $G_y$ satisfies $aI\preceq G_y\preceq I$ with a public unit
vector $v$. For $M$ inputs define
\begin{equation}\label{eq:direct-average}
 H_{\boldsymbol y}=\bigoplus_{j=1}^M G_{y^{(j)}},\qquad
 b=M^{-1/2}\bigoplus_{j=1}^M v,
 \qquad q(H_{\boldsymbol y},b)=M^{-1}\sum_jq(G_{y^{(j)}},v).
\end{equation}
If all inner instances have identical public norm tables and magnitude
patterns, global SQ first selects a uniform block and then uses its
inner SQ interface. Norms, supports and sparsity are unchanged within
a block. The same statement applies to sparse constraint factors:
their direct sum realizes \eqref{eq:direct-average} as one analytic-center
squared decrement, multiplied by the public factor $\eta^2/2$.

For exact levels $q_0,q_1=q_0(1+D)$ labelled by a partial Boolean
function $g$, an input with label-one fraction $p$ has scalar
$q_0(1+Dp)$. In the high-contrast range $D\geq1$, choosing
$p=\Theta(1/D)$ keeps the scalar of order $q_0$, while relative error
$\epsilon$ resolves a change $\Theta(\epsilon/D)$ in $p$.
Lemma~\ref{lem:weight} then gives an outer cost
$\Omega(\min\{M,D/\epsilon^2\})$; ordinary sampling supplies the
corresponding upper bound up to confidence factors. Thus a
$\kappa/\epsilon^2$ statistical factor from this averaging mechanism
requires contrast $1+D=\Theta(\kappa)$, rather than merely a large
absolute difference between two forms of order $\kappa$.

We call a realization \emph{oracle-preserving} if every granted sparse
count, location and value request, vector coordinate request, norm request,
and sampling request can be simulated with $O(1)$ queries to the hidden
input, with all remaining data public. This requirement also covers every
granted constraint-factor request. The constant is uniform over the
problem parameters and the requested operations.

\begin{proposition}[Conditional high-contrast compiler]\label{prop:contrast}
Let $g$ be a nonconstant partial Boolean problem with bounded-error randomized
query complexity $R(g)=L$. Suppose an oracle-preserving sparse SPD family
with identical public full-SQ metadata has the exact levels just described.
For $D\geq1$ and
sufficiently small $\epsilon$, one relative inverse-form estimate on
a direct sum of $M=\Theta(D/\epsilon^2)$ blocks requires
$\Omega(DL/\epsilon^2)$ queries.
\end{proposition}
\begin{proof}
Choose a small fixed $\gamma>0$, a weight
$t=\lfloor\gamma M/D\rfloor$, and
$\Delta=\lceil C_1\epsilon M/D\rceil$, with $C_1$ an absolute
constant large enough to separate relative output intervals.
Choose $M=C_2D/\epsilon^2$ rounded to an integer, and then choose
$\epsilon$ sufficiently small and $C_2$ sufficiently large.
These choices ensure $1\leq\Delta\leq t/4$, $t\leq M/4$,
$q_t=q_0(1+Dt/M)=\Theta(q_0)$, and
$q_{t+\Delta}-q_t\geq C_1\epsilon q_0$.
The outer partial function $F$ distinguishes weights $t,t+\Delta$.
Lemma~\ref{lem:weight} gives $R(F)=\Theta(M)$.
Theorem 2 of \citet{ChakrabortyEtAl2023} explicitly allows a partial
outer function and proves $R(F\circ g^M)=\Omega(R(F)R(g))$ whenever
$R(F)=\Theta(M)$. Equation~\eqref{eq:direct-average} and the disjoint
relative output intervals reduce this composed decision problem to
the scalar estimate with constant query overhead.
\end{proof}

If levels have relative half-width $\eta$, the weighted average can
move by $\eta q_0(1+Dp)$ even with all labels fixed. Thus the pointwise
reduction remains valid for $\eta\leq c\epsilon$ with a sufficiently
small absolute $c$. With one absolute half-width $\delta q_1$ for both
levels, it requires $\delta\leq c\epsilon/(1+D)$ instead. A constant-gap
Forrelation promise supplies neither of these narrow-level conditions.
An unrestricted assertion $R(F\circ g)=\Omega(R(F)R(g))$ is also
false for partial functions \citep{BenDavidBlais2020}; the linear
outer complexity in Proposition~\ref{prop:contrast} is material.

\subsection{A distributional product for a single positive form}
We recall precisely the distributional consequence of the composition
results that will be used. A noisy bit query of bias $\gamma$ returns
the correct bit with probability $(1+\gamma)/2$ and costs $\gamma^2$.
Write $\operatorname{noisyR}(F)$ for the minimum worst-case expected
cost of a bounded-error algorithm in this model. Theorems 24 and 35,
Definitions 33--34 and the minimax observation following Definition 34
of \citet{BenDavidBlais2020} imply the following: for every nonconstant partial
Boolean inner function $g$, there are distributions $\mu_0,\mu_1$ on
its two fibers such that, for every outer $F$, there is a distribution
$\nu$ on $\operatorname{Dom}(F)$ for which computing $F$ under
\begin{equation}\label{eq:product-hard-distribution}
 z\sim\nu,\qquad Y_j\sim\mu_{z_j}\ \text{independently conditional on }z
\end{equation}
requires $\Omega(\operatorname{noisyR}(F)R(g))$ expected queries.
Here the inner distributions can be chosen to witness the
Shaltiel-free complexity of $g$, which Theorem 24 lower-bounds by
$\Omega(R(g))$. This order of quantifiers is important: the same
$\mu_0,\mu_1$ work before $\nu$ is selected. The proof of their
Theorem 35 establishes this statement for that fixed pair; it is
stronger than simply maximizing the final composed complexity over
unstructured distributions. Worst-case query-bounded estimators are
covered by this expected-cost lower bound. We use the theorem
numbering of arXiv:2002.10809v2.

\begin{theorem}[Distributional scalar product]\label{thm:distributional-product}
Let $g$ be a nonconstant partial Boolean problem, and let its oracle-preserving
SPD realization have a public unit $v$ and
\begin{equation}\label{eq:single-form-signal}
 q(G_y,v)=d+C\Phi(y),\quad C>0,\quad d>0,\quad
 \chi=C/d\leq5/3,
\end{equation}
where $|\Phi|\leq1/100$ on zero inputs and $3/5\leq\Phi\leq1$ on
one inputs. Assume the identical public full-SQ metadata needed for
\eqref{eq:direct-average}. If $0<\epsilon\leq\chi/200$, there is a
family of $M=4\lceil\chi/\epsilon\rceil^2$ blocks for which relative
$\epsilon$ estimation of the \emph{single} averaged form requires
\begin{equation}\label{eq:distributional-product}
             \Omega((\chi/\epsilon)^2R(g))
\end{equation}
classical full-SQ queries. A factor realization transfers the same
statement to one box-LP squared decrement.
\end{theorem}
\begin{proof}
Set $n=\lceil\chi/\epsilon\rceil\geq200$, $M=4n^2$, and let $F$
distinguish the weights $M/2\pm24\sqrt M=2n^2\pm48n$.
They are integers in $[0,M]$. The Hamming-sphere argument proving
Lemma~\ref{lem:weight} applies equally for $t=\Theta(M)$ bounded away
from both endpoints: before a small fraction of $M$ queries, the two
Bernoulli probabilities are bounded away from zero and one, so the
per-query divergence is $O(1/M)$. Thus $R(F)=\Theta(M)$.
Observation 23 of \citet{ChakrabortyEtAl2023} gives
$\operatorname{noisyR}(F)=\Omega(R(F)^2/M)=\Omega(M)$.
The reverse bound follows by reading all bits.

Choose the distributions in \eqref{eq:product-hard-distribution}.
Let $\Delta=\E_{\mu_1}\Phi-\E_{\mu_0}\Phi\geq59/100$.
For each fixed outer input of one of the two weights, the conditional
mean of $M^{-1}\sum_j\Phi(Y_j)$ depends only on that weight.
Their means differ by $48\Delta/\sqrt M=24\Delta/n$.
Hoeffding's inequality, since every $\Phi$ lies in $[-1,1]$, gives
\[
 \Prob\left[\left|M^{-1}\sum_j\Phi(Y_j)
       -\E[M^{-1}\sum_j\Phi(Y_j)\mid z]\right|
       >\frac{24\Delta}{4\sqrt M}\ \middle|\ z\right]
 \leq2e^{-24^2\Delta^2/32}<0.004.
\]
Hence the two mean scalar values are separated by
$S=24\Delta C/n$, and each realized scalar is within $S/8$ of its
mean except on the displayed event. Because
$n\leq(201/200)\chi/\epsilon$, we have $S>14\epsilon d$.
Every scalar is at most $(1+\chi)d<3d$.
Amplify the estimator's success to $11/12$ by a constant number of
independent repetitions and a median. Its absolute error is then at
most $3\epsilon d$, and $S/8+3\epsilon d<S/2$.
Thresholding at the midpoint of the two public means computes $F$
with distributional error below $1/3$. This threshold may depend on
the fixed hard distributions; such nonuniform dependence is allowed
in the distributional lower bound. The composition theorem therefore
requires $\Omega(MR(g))$ queries. All scalar-input queries have
constant-cost hidden-input simulations, proving the result.
\end{proof}

For the clock of Theorem~\ref{thm:clock-lower}, gauge equivalence
makes the average of its two queried inverse diagonals a public number
$d$. Its positive polarization has exactly
$q_+(y)=d+\kappa c\Phi(y)$, where $c$ is the public normalized endpoint
coefficient. Thus $C=\kappa c$ and $\chi_c=\kappa c/d$.
We have $d\leq\kappa$, so $\chi_c\geq c$.
Positive definiteness bounds the absolute cross entry by the average
of its two diagonals; applying it to a high input with
$\Phi\geq3/5$ gives $\chi_c\leq5/3$.
The Montanaro--Shao clock at a sufficiently small coefficient scale $c$
has source hardness
\begin{equation}\label{eq:inner-Lc}
 L_c=\Omega\left(
 \frac{((s-1)/2)^{\Omega(\sqrt\kappa\log(1/c))}}
 {\log(s)\sqrt\kappa\log(1/c)}\right).
\end{equation}
One can use actual attainable coefficients within an absolute factor
of the desired scale, which changes no bounds. The explicit construction
below establishes this fact without assuming a continuum of clock sizes.
Theorem~\ref{thm:distributional-product} gives
$\Omega((\chi_c/\epsilon)^2L_c)$ for a single form whenever
$\chi_c\geq200\epsilon$. In particular a fixed small coefficient gives
$\epsilon^{-2}s^{\Omega(\sqrt\kappa)}$ up to the displayed denominator,
on one instance. For coefficient scales $c=c_0\epsilon^\alpha$,
$0\leq\alpha<1$, the conservative $\chi_c\geq c$ gives
\[
 \epsilon^{-2(1-\alpha)}
 s^{\Omega(\sqrt\kappa[1+\alpha\log(1/\epsilon)])},
\]
again with the denominator of \eqref{eq:inner-Lc} and sufficiently
small $\epsilon$. This is a simultaneous-cost statement; it does
not establish a product of both full factors in the upper bound.

\subsection{A witness clock with an exact diagonal}\label{sec:witness-diagonal}
The following calculation determines the signal-to-baseline ratio for
an explicit clock. It uses a restricted interval to produce a valid
full-domain dual witness, rather than asserting that the restricted
approximant is minimax on all of $[-1,1]$.

\begin{proposition}[Explicit witness diagonal]\label{prop:witness-diagonal}
For every $\kappa\geq4$ and integer $r\geq1$, there is an odd-order
zero-diagonal Jacobi contraction $J$ of order $2r+5$ with positive
edges, whose second and penultimate coordinates satisfy
\begin{equation}\label{eq:witness-values}
 (f_\kappa(J))_{2,2r+4}=c_r,\qquad
 (f_\kappa(J))_{2,2}=(f_\kappa(J))_{2r+4,2r+4}=d',
\end{equation}
where
\[
 c_r=\frac{C\delta^r}{A^2-1},\qquad
 d'=\frac{CA}{A^2-1}
 =\frac{(\kappa+1)(\kappa^2+14\kappa+1)}
 {4\kappa(\kappa^2+6\kappa+1)},
\]
with $A=4(m/h)^2-3$, $C=4am/h^2$, and
$\delta=A-\sqrt{A^2-1}$. In particular $1/4<d'<3/5$,
$\chi_{c_r}=c_r/d'=\Theta(c_r)$, and
$r=\Theta(\sqrt\kappa\log(1/c_r))$ at sufficiently small scales.
\end{proposition}
\begin{proof}
A full derivation is given here because the diagonal, unlike the
approximation degree alone, controls statistical composition.
The even part of $f_\kappa$ is $am/(m^2-h^2x^2)$.
On $x^2\in[1/2,1]$, set $t=4x^2-3$; it becomes $q(t)=C/(A-t)$.
Put $\omega=\sqrt{A^2-1}$. The degree-$r$ reciprocal minimax error is
$E_r=C\delta^r/\omega^2$. Write $T_j$ and $U_j$ for the first- and
second-kind Chebyshev polynomials, respectively:
$T_j(\cos\varphi)=\cos(j\varphi)$ and
$U_j(\cos\varphi)=\sin((j+1)\varphi)/\sin\varphi$, with continuous
extension at the endpoints. The scalar polynomial $U_j$ is distinguished
from the matrix-valued clock gate $U_t$ by its argument and context.
The normalized error is
\[
 V_r(t)=\frac{N_r(t)}{A-t},\qquad
 N_r(t)=(At-1)T_r(t)+\omega(t^2-1)U_{r-1}(t).
\]
To verify rather than assume this expression, write $t=\cos\varphi$.
The half-angle identity
$\tan(\psi/2)=\sqrt{(A+1)/(A-1)}\tan(\varphi/2)$ gives
$V_r(t)=\cos(r\varphi+\psi)$. Also
$N_r(A)=\omega^2(A+\omega)^r$, so $C-E_rN_r(t)$ vanishes at
$t=A$; dividing by $A-t$ yields a degree-$r$ polynomial.
The phase increases strictly from zero to $(r+1)\pi$, producing
$r+2$ alternating extremal points. The alternation theorem therefore
verifies the minimax claim.

Let these points be $t_0<\cdots<t_{r+1}$, and define
\[
 R_r(t)=(At-1)U_{r-1}(t)+\omega T_r(t),\qquad
 L_r(t)=(t^2-1)R_r(t).
\]
The interior extrema are exactly the $r$ roots of $R_r$ (use
$\sin(r\varphi+\psi)=0$). Thus the $t_j$ are the simple roots of
$L_r$, whose leading coefficient is positive, and
$V_r(t_j)=\operatorname{sgn}L_r'(t_j)=:\sigma_j$.
Set $\mu_j=|L_r'(t_j)|^{-1}/\sum_k|L_r'(t_k)|^{-1}$.
Partial fractions at infinity show
\begin{equation}\label{eq:barycentric-moments}
 \sum_j\frac{P(t_j)}{L_r'(t_j)}=0\quad(\deg P\leq r).
\end{equation}
Hence the signed weights $\sigma_j\mu_j$ annihilate every
polynomial of degree at most $r$ and pair with $q$ to $E_r$.
At $x_j=\sqrt{(t_j+3)/4}$, put signed mass $\sigma_j\mu_j/2$
on each of $\pm x_j$. Odd moments cancel, and even moments through
degree $2r$ vanish by \eqref{eq:barycentric-moments}. This is a
full-domain dual witness supported inside $[-1,1]$.

We now realize its cross entry explicitly, in the form underlying
\citet[Lemma 2.5]{MontanaroShao2023}. Put
$B=(2\sum_j\mu_j/x_j^2)^{-1}$ and choose the positive symmetric
probability measure with mass $1/2$ at zero and
$B\mu_j/(2x_j^2)$ at each $\pm x_j$.
Gram--Schmidt on $1,x,\ldots,x^{2r+4}$ gives its Jacobi matrix $J$.
The symmetry makes its diagonal zero, positivity of the distinct-node
measure makes all edges positive, and its eigenvalues are the support
points, so it is a contraction. Since $\int x^2\,dw=B$, the second
orthonormal polynomial is $x/\sqrt B$; its spectral masses at
$\pm x_j$ are $\mu_j/2$ and at zero are zero.
Interpolate the values $\sigma_j$ by a degree-$(r+1)$ polynomial
$S(t)$. Equation~\eqref{eq:barycentric-moments} implies that
$xS(4x^2-3)/\sqrt B$ is orthogonal to every lower odd polynomial
and, by symmetry, every even polynomial. Its norm is one, since
$S(t_j)^2=1$. Its leading coefficient is positive: the divided
difference of the ordered alternating values is a sum of positive
terms. Thus it is exactly the penultimate orthonormal polynomial,
of degree $2r+3$. Its spectral masses agree with those of the second
coordinate, and their signed cross masses are $\sigma_j\mu_j/2$.
The cross entry is consequently $E_r=c_r$, and both diagonals are
$d'=\sum_j\mu_jq(t_j)$.

To evaluate that average, partial fractions of $N_r/L_r$ give
\[
 S_0:=\sum_j|L_r'(t_j)|^{-1}=\frac{N_r(A)}{L_r(A)},\qquad
 S_1:=\sum_j\frac{q(t_j)}{|L_r'(t_j)|}
       =-C\left(\frac{N_r}{L_r}\right)'(A).
\]
Indeed $N_r(t_j)/(A-t_j)=\sigma_j$, and differentiating the
partial-fraction expansion gives the second identity. If
$\mathcal S=T_r(A)+\omega U_{r-1}(A)=(A+\omega)^r$, then
$N_r(A)=\omega^2\mathcal S$ and $L_r(A)=\omega^3\mathcal S$.
Using $T_r'=rU_{r-1}$ and
$\omega^2U_{r-1}'(A)=rT_r(A)-AU_{r-1}(A)$ gives
\[
 \frac{N_r'(A)}{N_r(A)}=\frac r\omega+\frac A{\omega^2},\qquad
 \frac{L_r'(A)}{L_r(A)}=\frac r\omega+\frac{2A}{\omega^2}.
\]
It follows that $d'=S_1/S_0=CA/(A^2-1)$, as claimed.
The rational expression and $1/4<d'<3/5$ follow by substitution
and multiplication of positive denominators. Finally
$\log(1/\delta)=\operatorname{arcosh}(A)=\Theta(\kappa^{-1/2})$
and $C/(A^2-1)=\Theta(1)$ for $\kappa\geq4$.
\end{proof}

In the lifted clock use public identity gates before and after the
hard circuit, so its second and penultimate positions carry the
Forrelation amplitude. Fitting an unpadded depth-$\Theta(r)$ circuit
between these positions gives the exponent in \eqref{eq:inner-Lc};
a bounded number of identity gates adjusts the residue of the length.
Successive coefficients $c_r$ have ratio $\delta$, bounded below by
an absolute constant for $\kappa\geq4$, so there is always an
attainable coefficient within a constant factor of any sufficiently
small requested scale. The inverse baseline is $d=\kappa d'$.
Consequently this explicit family has statistical factor
$\Theta((c_r/\epsilon)^2)$, without a hidden condition-dependent gain.

\subsection{A constant-weight path at high accuracy}
Another clock gives a smaller baseline and a different product.
\begin{theorem}[High-accuracy path product]\label{thm:path-product}
There are absolute constants $c_*>0$ and $\kappa_0$ such that for
$\kappa\geq\kappa_0$, $0<\epsilon\leq c_*\kappa^{-3/2}$, and
$q=2^r$ with integer $r\geq1$, a full-SQ family with matrix sparsity
at most $2q+1$, condition exactly $\kappa$, and one public right-hand
side requires
\begin{equation}\label{eq:path-product}
 \Omega\left(\frac{\kappa^2q^{\ell/2}}{r\ell}\right),\qquad
 \ell=\Theta\left(\sqrt\kappa
        \left[1+\log\frac1{\kappa^{3/2}\epsilon}\right]\right)
\end{equation}
queries to estimate its inverse form to relative error $\epsilon$. It has an
$O(q)$-sparse box-LP realization with the same query simulation.
\end{theorem}
\begin{proof}
Take an $n$-vertex scalar path $J_n$ with all off-diagonals $1/2$.
Lift its transitions to the Forrelation circuit as in
\eqref{eq:weighted-clock}. Let $H=mI+hJ_y$, and append the scalar
pads $a,1$. Set
$z_0=m/h=(\kappa+1)/(\kappa-1)$,
$R=z_0+\sqrt{z_0^2-1}=(\sqrt\kappa+1)/(\sqrt\kappa-1)$,
and $\theta=\log R$.
The determinant recurrence, with $D_0=1,D_1=m$, gives
$D_n=(h/2)^nU_n(z_0)$.
Cofactors give the normalized endpoint coefficient and diagonal
\begin{equation}\label{eq:path-values}
 c_n=a|(H^{-1})_{1n}|=\frac{2a}{hU_n(z_0)},\qquad
 d_n'=a(H^{-1})_{11}=\frac{2aU_{n-1}(z_0)}{hU_n(z_0)},\qquad
 \chi_n=\frac{c_n}{d_n'}=\frac1{U_{n-1}(z_0)}.
\end{equation}
The opposite endpoint has the same diagonal. Choose the sign of the
public two-sparse vector to cancel the alternating cofactor sign;
its form is $\kappa d_n'+\kappa c_n\Phi$.
Since $U_j(z_0)=\sinh((j+1)\theta)/\sinh\theta$, we obtain
\[
 c_n=\frac{16\sqrt\kappa}{(\kappa-1)^2}
       \frac1{R^{n+1}-R^{-(n+1)}},\quad
 \chi_n=\frac{4\sqrt\kappa}{\kappa-1}
       \frac1{R^n-R^{-n}},\quad
 \frac{\chi_n}{c_n}=\frac{\kappa-1}{4}
       \frac{\sinh((n+1)\theta)}{\sinh(n\theta)}.
\]
Uniformly for $n\theta\geq1$ and $\kappa\geq16$ this implies
$c_n=\Theta(\kappa^{-3/2}e^{-n\theta})$,
$d_n'=\Theta(\kappa^{-1})$ and $\chi_n=\Theta(\kappa c_n)$.

To make increasing length genuine source hardness, take a work register
of $r\ell$ qubits, decompose each of the three full Hadamards into
$\ell$ disjoint $r$-qubit layers, and use the two hidden sign queries.
There are exactly $3\ell+2$ transitions and
$n=3(\ell+1)$ vertices. Choose the largest such $n$ with
$c_n\geq\epsilon$. For sufficiently small $c_*$ this length has
$n\theta\geq1$. Consecutive admissible lengths differ by three and
the displayed formula bounds their coefficient ratio by an absolute
constant. Hence $c_n=\Theta(\epsilon)$, and solving the exponential
formula gives the stated $\ell$. Further,
$\chi_n\geq(\kappa-1)\epsilon/4$.
Choose, for example, $\kappa_0=1024$ so
$\chi_n\geq200\epsilon$. The source Forrelation problem has
randomized complexity $\Omega(2^{r\ell/2}/(r\ell))$.
Apply Theorem~\ref{thm:distributional-product} to
$M=4\lceil\chi_n/\epsilon\rceil^2=\Theta(\kappa^2)$ blocks.
This proves \eqref{eq:path-product}.

The unlifted path has norm less than one, so the affine shift obeys
the promised spectral bounds, and the public pads make the condition
exact. Every transition has public squared row and column mass $1/4$.
Interior squared row norms of $H$ are $m^2+h^2/2$ and boundary
norms are $m^2+h^2/4$. All conditional and global samples are public
mixtures of uniform Hadamard supports, a single sign location, and
the diagonal; values cost one hidden query. Sparsity is at most
$2q+1$. Proposition~\ref{prop:box-factor} supplies the sparse factor
with the same simulation. Direct sums preserve these properties.
\end{proof}

The larger constant $\kappa_0$ is only convenient bookkeeping for the
concentration theorem. This refinement is restricted to high accuracy:
the constant path's boundary prefactor $\Theta(\kappa^{-3/2})$
shortens the admissible hard-circuit length, relative to the optimized
witness at the same coefficient scale, by order
$\sqrt\kappa\log\kappa$. Its explicit $\kappa^2$ outer factor
therefore does not imply an improved leading exponential envelope.

\subsection{Proved limitations of particular contrast constructions}
\begin{proposition}[Local cyclic-clock dilution]\label{prop:cyclic-dilution}
Let $U$ be unitary on an $L$-position clock, and let $e$ satisfy
$U^Le=\sigma e$ for $\sigma\in\{1,-1\}$, with
$e,Ue,\ldots,U^{L-1}e$ orthonormal. For $0<\gamma<1$ set
$A=I-\gamma U$, $G=AA^*/(1+\gamma)^2$, and
$K=(1+\gamma)/(1-\gamma)$. Then $\kappa(G)\leq K^2$ and
\begin{equation}\label{eq:cyclic-contrast}
 \frac{q_+}{q_-}
 =\left(\frac{1+\gamma^L}{1-\gamma^L}\right)^2
 \leq(1+K/L)^2,
 \qquad q_\sigma=e^*G^{-1}e.
\end{equation}
\end{proposition}
\begin{proof}
The singular values of $I-\gamma U$ lie between $1-\gamma$ and
$1+\gamma$. The finite geometric identity gives
$A^{-1}e=(1-\sigma\gamma^L)^{-1}\sum_{t=0}^{L-1}\gamma^tU^te$.
Since $A$ is normal, $q_\sigma=(1+\gamma)^2\norm{A^{-1}e}^2$.
Orthogonality makes the numerator $\sum_t\gamma^{2t}$ common to
both signs, proving the equality. Write the ratio as
$\coth^2((L/2)\log((K+1)/(K-1)))$. The inequalities
$\log((K+1)/(K-1))\geq2/K$ and
$\coth x\leq1+1/x$ prove the upper bound.
\end{proof}
For a real circuit this is an ordinary real SPD construction. If each
transition is $q$-sparse in rows and columns, $G$ is $(2q+1)$-sparse.
Public scalar pads can make its condition exactly $K^2$.
Equation~\eqref{eq:cyclic-contrast} bounds contrast by
$O(1+K^2/L^2)$ even for an \emph{exact} holonomy sign. Thus contrast
of order the final condition requires bounded $L$ in this construction.
This is a limitation of the stated first-order cyclic resolvent, not of
all sparse linearizations.

\begin{proposition}[Two-cluster and Schur-complement limitations]\label{prop:contrast-obstructions}
The following two assertions hold.
\begin{enumerate}
\item If $J_y=J_y^T$, $J_y^2=I$, and orthogonal public unit vectors
$x,z$ satisfy $x^TJ_yx=z^TJ_yz=0$, put $u_y=x^TJ_yz$ and
$G_y=mI+hJ_y$. For $v=(x+z)/\sqrt2$,
\begin{equation}\label{eq:two-cluster}
 q(G_y,v)=1+\frac{\kappa-1}{2}(1-u_y).
\end{equation}
A high-amplitude interval $u_y\in[3/5,1]$ therefore fails to force
an $O(1)$ scalar uniformly as $\kappa$ grows. Such a uniform bound
requires $1-u_y=O(1/\kappa)$.
\item Let $G\succ0$ have largest eigenvalue one and condition $K_0$.
For an SPD extension
$\mathcal H=\left(\begin{smallmatrix}\alpha&w^T\\w&G\end{smallmatrix}\right)$
with Schur complement $S=\alpha-w^TG^{-1}w>0$,
\begin{equation}\label{eq:schur-condition}
 e_0^T\mathcal H^{-1}e_0=S^{-1},\qquad
 \kappa(\mathcal H)\geq\frac{1+\norm{G^{-1}w}^2}{S}.
\end{equation}
If $w=tv$ with $\norm v=1$, $v^TG^{-1}v\geq cK_0$,
$w^TG^{-1}w\geq c_0>0$, and $S=\Theta(1/R)$, then
$\kappa(\mathcal H)=\Omega(K_0R)$.
\end{enumerate}
\end{proposition}
\begin{proof}
Since $m^2-h^2=a$, $(mI+hJ)^{-1}=(mI-hJ)/a$.
Substitution gives \eqref{eq:two-cluster}. At $u=3/5$ its value is
$\kappa/5+4/5$, while for $|u|\leq1/100$ it is of order $\kappa$;
the uniform contrast between those intervals is only constant.
For the second assertion, block inversion gives the inverse entry.
The vector $(1,-G^{-1}w)$ has Rayleigh quotient
$S/(1+\norm{G^{-1}w}^2)$, and interlacing gives
$\lambda_{\max}(\mathcal H)\geq1$, proving
\eqref{eq:schur-condition}. Finally Cauchy--Schwarz for the spectral
measure of $v$ gives $v^TG^{-2}v\geq(v^TG^{-1}v)^2$. Consequently
\[
 \norm{G^{-1}w}^2=t^2v^TG^{-2}v
 \geq(w^TG^{-1}w)(v^TG^{-1}v)\geq c_0cK_0.
\]
Combine this with \eqref{eq:schur-condition} and $S=\Theta(1/R)$.
\end{proof}

Under the hypotheses of the Schur construction, a final condition budget
$\overline K$ enforces $R=O(\overline K/K_0)$. If its formal clock
factor is $\exp(A\sqrt{K_0})$ for fixed $A>0$, optimizing the
condition-dependent part of the proposed product gives
\[
 \max_{1\leq K_0\leq\overline K}
 \frac{\overline K}{K_0}e^{A\sqrt{K_0}}
 =\max\{\overline K e^A,e^{A\sqrt{\overline K}}\}.
\]
Indeed the logarithm, in $t=\sqrt{K_0}$, is
$\log\overline K-2\log t+At$, which is convex and hence maximized
at an endpoint. This algebra explains the condition-budget loss for
this specific compiler. It is not a general upper bound for sparse
SPD query complexity, and it does not rule out the common
$\epsilon^{-2}$ factor already proved distributionally.

There is also a simple circuit-resource qualification. If a chosen
coherent repetition or polynomial amplification procedure uses $r_{\rm amp}$ invocations of a depth-$D$ source circuit, placing that
procedure in a clock of length $T$ requires
$r_{\rm amp}D\leq O(T)$. Standard repetition to reduce a constant
decision error to $O(\epsilon/\kappa)$ uses
$r_{\rm amp}=O(\log(\kappa/\epsilon))$ invocations.
One cannot count the original depth-$T$ Forrelation exponent and then
amplify its promise at no clock-length cost. This bookkeeping does not
assert that every structure-specific amplification must have that cost;
nor does a narrow acceptance probability by itself supply the
condition-preserving positive-form realization in
Proposition~\ref{prop:contrast}.

Our conclusions are therefore separate, explicit lower families:
\eqref{eq:clock-lower}, \eqref{eq:statistical-lower},
\eqref{eq:distributional-product}, and \eqref{eq:path-product}.
For worst-case complexity over the corresponding dimensions their
maximum is a valid lower envelope. We do not claim the full product
$\kappa\epsilon^{-2}s^{\Omega(\sqrt\kappa\log(1/\epsilon))}$.
It would require additional information, for example a hard single-form
family with condition-scale contrast and the narrow levels of
Proposition~\ref{prop:contrast}. The proved products and the exact
construction limitations stand independently of that stronger question.
